\subsection{Two reflection maps reduce genus two to the ordinary case}
\label{sec:quintic-ordinary-reduction}

\paragraph{Finding A33 (a quintic transfer to an ordinary genus-two curve).}
For a curve over a finite field of characteristic two, let $\Pi_C$ have
the meaning of A1: the geometric kernel is replaced by its maximal
prime-to-two quotient, and the arithmetic quotient remains the whole
absolute Galois group.  Consider the following restricted assertion:
\[
 \begin{gathered}
 \text{For every finite field $k$ of characteristic two and every ordinary}\\
 \text{smooth proper geometrically connected genus-two curve $E/k$,}\\
 E(k)\longrightarrow\operatorname{Sec}(\Pi_E\to G_k)/\Delta_E
 \text{ is injective.}
 \end{gathered}
 \tag{DR.1}
\]
If (DR.1) holds, the same point-section assertion holds for every
smooth proper geometrically connected genus-two curve over every finite
field of characteristic two, without a restriction on its $2$-rank.
In fact the hypothesis can be restricted further to ordinary $E/k$
equipped with an embedding
$\mathcal O_5=\mathbf Z[(1+\sqrt5)/2]\hookrightarrow
\operatorname{End}_k(J_E)$ whose image is fixed by the Rosati involution
of the canonical principal polarization.  Thus only this real
multiplication case of the ordinary assertion is needed.

More precisely, let $C/\overline k$ have genus two and let
$\pi:D\to C$ be any connected cyclic etale cover of degree five, with
deck generator $\sigma$.  A lift $\tau$ of the hyperelliptic involution
can be chosen, and it satisfies
\[
 \tau^2=1,\qquad \tau\sigma\tau^{-1}=\sigma^{-1}.
 \tag{DR.2}
\]
The quotient $v:D\to E=D/\langle\tau\rangle$ has genus two.  If
$\mathcal B=\ker(\operatorname{Nm}_\pi:J_D\to J_C)^0$, then
\[
 \dim\mathcal B=4,\qquad
 \mathcal B\sim J_E\times J_E,\qquad
 f_{\mathcal B}\in\{0,4\},\qquad f_E\in\{0,2\}.
 \tag{DR.3}
\]
Here $\sim$ is isogeny over the algebraic closure and $f$ denotes
geometric $2$-rank.  After a finite constant extension all the maps and
deck transformations can be defined and a chosen base-point fibre can
be split.  If $f_{\mathcal B}=0$, this cover itself supplies the
geometric primitive-zero witness.  If $f_{\mathcal B}=4$, the two maps
\[
                         v,\quad v\circ\sigma:D\longrightarrow E
 \tag{DR.4}
\]
transfer any one actual point-section collision on $C$ to two collisions
on the same ordinary genus-two curve $E$.  Their fibres together
separate geometric points of $D$.

\paragraph{Source attribution and scope.}
The reflection-lift, completed-local-ring, and Prym-image arguments
already occur in recovered A1 B.40 for cubic covers in genera two and
three \cite[B.40]{A1run}.  Recovered successor A.37 studies a pointed
cubic reflection quotient and explains why a reflection twist does not
remove a two-primary obstruction on its elliptic quotient
\cite[A.37]{ASrun}.  The complete proof of successor A.4 explicitly
lifts equal point sections to equal point sections on the same chosen
odd cover \cite[A.4]{ASrun}.  We use that machinery explicitly here.  The
degree-five choice retains genus two, and the two quotient maps below
transport one common lifted section collision.  This is a reduction of
the older route to its ordinary genus-two case; it supplies neither a
universal theta-incidence assertion nor a proof of (DR.1) by Hooley.
No analytic hypothesis is used in the reduction.

\begin{proof}
\emph{1. The reflection quotient and its genus.}
Work first over an algebraic closure.  Connected cyclic etale covers of
degree five exist because $J_C[5]\simeq(\mathbf Z/5)^4$ geometrically.
One may use an exact-order-five line bundle and a trivialization of its
fifth power.  For any such cover, the hyperelliptic involution $\iota$
acts on $J_C$ by $[-1]$, and hence takes its defining torsion line
bundle to the inverse.  It therefore lifts to an automorphism $\tau$
of $D$ that inverts the deck group.  Its square is some $\sigma^a$.
Conjugating that square by $\tau$ both leaves it unchanged and replaces
it by $\sigma^{-a}$.  Thus $2a=0$ in $\mathbf Z/5$, so $a=0$ and
(DR.2) follows.

The hyperelliptic map $C\to\mathbf P^1$ is a separable double cover in
characteristic two.  Above each of its geometric branch points,
$\iota$ has one fixed point $W$ on $C$.  The five points of
$\pi^{-1}(W)$ form a deck torsor, on which $\tau$ acts as
$j\mapsto b-j$.  This has exactly one fixed point, because two is
invertible modulo five.  Every fixed point of $\tau$ occurs this way,
since its image is fixed by $\iota$.

At each such fixed point $T$, the etale map $\pi$ identifies the
completed local rings at $T$ and $W$.  The actions of $\tau$ and
$\iota$ agree under that identification.  Consequently the local
different exponent of $D\to E$ at $T$ is exactly the different
exponent of $C\to\mathbf P^1$ at $W$.  There is no ramification away
from these fixed points.  The total different of the latter
hyperelliptic map is six by Riemann--Hurwitz.  Since $g(D)=6$ by the
unramified degree-five formula, Riemann--Hurwitz for $D\to E$ gives
\[
                 10=2g(D)-2=2(2g(E)-2)+6,
                 \qquad g(E)=2.
 \tag{DR.5}
\]
This argument works for all three possible $2$-ranks of $C$; it does
not assume the rank-one branch pattern.

\emph{2. The Prym isogeny and the rank dichotomy.}
Let $A=v^*J_E\subset J_D$ and let $A'=\sigma A$.  Pullback has
finite kernel because its composition with the corresponding norm is
$[2]$, so both images have dimension two.  The norm of $A$ to $J_C$
is fixed by the induced hyperelliptic involution, which acts by $[-1]$.
The homomorphism $\operatorname{Nm}_\pi\circ v^*$ is therefore killed
by two.  A homomorphism of abelian varieties killed by a nonzero integer
is zero, so $A\subset\mathcal B$; the same is true of $A'$.

Every point of $A$ is fixed by $\tau$, and every point of $A'$ is fixed
by $\sigma\tau\sigma^{-1}$.  These two reflections generate the
degree-five rotation group.  Thus a point in their intersection is
fixed by $\sigma$.  On the other hand,
\[
 \pi^*\operatorname{Nm}_\pi=1+\sigma+\sigma^2+\sigma^3+\sigma^4.
 \tag{DR.6}
\]
Its norm is zero, so (DR.6) shows that the intersection is killed by
five.  This also proves scheme-theoretic finiteness: the intersection
is contained in $J_D[5]$.  The sum of $A$ and $A'$ therefore has
dimension four and equals $\mathcal B$.  Thus the map
$(x,y)\mapsto v^*x+\sigma v^*y$ from $J_E\times J_E$ to
$\mathcal B$ is an isogeny, as asserted in (DR.3).

There is a short group-algebra proof of the rank restriction.  The
geometric group $\mathcal B[2](\overline k)$ is an
$\mathbf F_2[\mathbf Z/5]$-module of dimension $f_{\mathcal B}\le4$.
It has no nonzero invariant vector: such a vector $x$ would satisfy
$0=(1+\sigma+\cdots+\sigma^4)x=5x=x$ by (DR.6).
The group algebra is semisimple, and its only nontrivial simple module
has dimension four because
\[
             \operatorname{ord}_5(2)=4,
             \qquad \Phi_5(T)=T^4+T^3+T^2+T+1
             \text{ is irreducible over }\mathbf F_2.
 \tag{DR.7}
\]
It follows that $f_{\mathcal B}$ is zero or four.  Isogeny invariance
and additivity of $2$-rank give $f_{\mathcal B}=2f_E$, proving all of
(DR.3).  In particular a nonzero new rank forces $E$ to be ordinary.

The quotient Jacobian also has the asserted integral real
multiplication.  Define the integral correspondence endomorphism
\[
                   T=v_*\sigma v^*\in\operatorname{End}(J_E).
 \tag{DR.7a}
\]
The pullback--norm identity $v^*v_*=1+\tau$ and the relation
$\tau\sigma=\sigma^{-1}\tau$ give
\[
            v^*T=(\sigma+\sigma^{-1})v^*.
 \tag{DR.7b}
\]
On $\mathcal B$, the deck sum in (DR.6) vanishes.  Expanding
$(\sigma+\sigma^{-1})^2+(\sigma+\sigma^{-1})-1$ gives exactly
$1+\sigma+\sigma^2+\sigma^3+\sigma^4$.  Therefore
$v^*(T^2+T-1)=0$.  Since $v^*$ has finite kernel, the homomorphism
$T^2+T-1$ has finite image and hence is zero.  The polynomial
$X^2+X-1$ is irreducible over $\mathbf Q$.  The induced unital map
from its quadratic integer ring to $\operatorname{End}(J_E)$ is
injective: any nonzero ideal in that ring contains a nonzero integer,
whereas multiplication by a nonzero integer on $J_E$ is not zero.
This is the desired copy of $\mathcal O_5$.

Under the canonical principal polarizations, norm and pullback are
adjoint and the adjoint of a curve automorphism is its inverse.
Consequently
\[
 T^{\dagger}=v_*\sigma^{-1}v^*
            =v_*\tau\sigma\tau v^*=v_*\sigma v^*=T.
 \tag{DR.7c}
\]
Here $v_*\tau=v_*$ and $\tau v^*=v^*$ were used.  The entire
quadratic ring is therefore fixed by Rosati.  This correspondence is
defined over any field over which $v$ and $\sigma$ are defined.

\emph{3. One common lift of a section collision.}
Let $P,Q\in C(k)$ give conjugate sections in the precise quotient
$\Pi_C$ used by A1.  Choose one geometric degree-five cover and the
reflection data above.  Since all coefficients and geometric points
are algebraic over a finite field, a finite constant extension $k'/k$
makes the cover, its deck transformations, and $\tau$ defined and makes
some chosen $P'\in D(k')$ over $P$ rational.  The constant degree-five
deck group then makes the entire fibre above $P$ split.  The assumed
section conjugacy persists on restriction to $G_{k'}$.

The arithmetic group of this pointed cover is an open subgroup of
$\Pi_{C_{k'}}$, and its geometric kernel has the full prime-to-two
topology.  Here is the normal-core justification needed for that
last assertion.  For the normal index-five subgroup defining this
cover, take the kernel of any additional finite odd quotient and
intersect its five conjugates in the original geometric fundamental
group.  The resulting quotient of the index-five subgroup embeds in
a product of finite odd groups.  Its extension by the original cyclic
deck group of order five again has odd order.  Thus that quotient
already factors through the original maximal pro-odd geometric
quotient.  This proves the compatibility with $\Pi_D$ used here and
in A01; the arithmetic Galois quotient has not been replaced.

Equivalently, represent the assumed geometric section conjugacy by a
$G_{k'}$-fixed prime-to-two path from $P$ to $Q$.  Its unique lift with
initial point $P'$ has a rational endpoint $Q'\in D(k')$ over $Q$ and
is a $G_{k'}$-fixed path between $P'$ and $Q'$.  Hence
\[
             [s_{P'}]=[s_{Q'}]
             \quad\text{in }\operatorname{Sec}(\Pi_D\to G_{k'})/\Delta_D.
 \tag{DR.8}
\]
This is the full-section lift checked in recovered successor A.4
\cite[A.4]{ASrun}.  There is one fixed pair of lifts, chosen before
either quotient map is applied.  It is not an inference from two
separately selected abelian Kummer lifts.

\emph{4. The zero-rank alternative already separates the points.}
Suppose $P\ne Q$ and $f_{\mathcal B}=0$.  Abelian Kummer theory applied
to (DR.8), as checked in A01, gives
\[
                     \delta=[Q'-P']\in J_D(k')[2^\infty].
 \tag{DR.9}
\]
The connected image of $\sigma-1$ is $\mathcal B$.  It lies there by
the norm identity, and conversely $5-(1+T+\cdots+T^4)$ is divisible
by $T-1$; on $\mathcal B$, substitution $T=\sigma$ expresses $[5]$
as a multiple of $\sigma-1$.  Surjectivity of $[5]$ on
$\mathcal B$ proves the reverse inclusion.  Since this abelian variety
has no nonzero geometric two-primary point, (DR.9) gives
$(\sigma-1)\delta=0$.

It follows that $\sigma Q'+P'-Q'-\sigma P'$ is principal.  Its four
points are distinct: the two base fibres are disjoint and the deck
action is free.  A function with this divisor has degree two and
simple poles, so it is separable in characteristic two.  Thus $D$ is
hyperelliptic.  Its genus is six, so its hyperelliptic pencil is unique.
The order-five deck automorphism descends to an order-five
automorphism of that projective line, because the kernel of the
descent action has order at most two.  Over the algebraic closure it
has a fixed point.  Its fibre has at most two geometric points, and a
group of order five acts trivially on such a set.  This gives a fixed
point of $\sigma$ on $D$, contrary to etaleness of the cyclic cover.
This is the primitive-zero deck argument of A11/A18, with its complete
degree-five specialization included here.

\emph{5. The ordinary alternative and the two maps.}
It remains to consider $f_{\mathcal B}=4$.  By Step 2, $E/k'$ is an
ordinary genus-two curve with the specified Rosati-invariant
$\mathcal O_5$ action defined over $k'$.  Every morphism of curves induces a map on
the arithmetic fundamental groups used here.  Indeed a pullback of
an etale cover is etale, regardless of whether the morphism itself is
ramified, and the pullback of a geometric odd quotient still has odd
monodromy.  Thus both maps in (DR.4) take (DR.8) to equal point-section
classes on $E$.

Assume (DR.1), or only its stated real multiplication restriction.
Applying it first to $v$ and then to $v\circ\sigma$
gives
\[
                 v(P')=v(Q'),\qquad
                 v(\sigma P')=v(\sigma Q').
 \tag{DR.10}
\]
The geometric fibres of $v$ are exactly the $\tau$-orbits.  If
$Q'\ne P'$, the first equality gives $Q'=\tau P'$.  The second then
gives $Q'=\sigma^{-1}\tau\sigma P'=\sigma^{-2}\tau P'$.
Consequently $\sigma^2$ fixes $\tau P'$, impossible for the free
degree-five deck action.  Therefore $P'=Q'$, and hence $P=Q$.
This proves the asserted reduction.  The ordinary assertion is also
plainly a special case of the unrestricted genus-two assertion, so
these two universal finite-field assertions are equivalent.
\end{proof}

\paragraph{A finite canonical-class sufficient condition after transfer.}
The ordinary universal assertion in (DR.1) can be replaced, for a
particular prescribed $P$ and a particular cover in the ordinary
alternative, by the two individual hypotheses of A29 at
\[
                 R_0=v(P'),\qquad R_1=v(\sigma P').
 \tag{DR.11}
\]
Namely, for each $i=0,1$ suppose the class
$\kappa_{R_i}=[K_E-2R_i]\in J_E(k')$ has a nonzero $r_i$-primary
component for some odd prime $r_i$ with even
$\operatorname{ord}_{r_i}(2)$.  The primes may differ.  A29 separates
the section of each $R_i$ from every other rational point of $E$.
Applying that result to the two pushed collisions proves (DR.10), so
the same proof separates $P$ on $C$.

If $P$ is Weierstrass, choose the reflection that fixes $P'$.  Its
quotient fibre over $R_0$ has the single geometric point $P'$, so the
one hypothesis at $R_0$ suffices; no second detector is then needed.
Such a reflection exists because the five reflections act on the
five-point fibre as all maps $j\mapsto b-j$.  This gives a legitimate
odd-cover transfer even when $\kappa_P=0$ on the original curve.
An A29 cover of $E$ pulls back to an etale odd cover of $D$; after
taking the pointed connected component and an odd normal closure if
needed, its composition with $D\to C$ remains geometrically of odd
degree.  The ramified double quotient is used only as a morphism
through which to pull back such covers, not as an allowed etale
double cover in the original prime-to-two tower.

The displayed canonical-class hypotheses are additional finite-group
conditions.  The reflection construction does not prove them.  Nor
does it force a theta point on $C$ of order five or of any other
specified order.  Thus the reduction avoids the genus-three
truncation obstruction of A30, but leaves an ordinary genus-two
existence question in the older route.
